\section{Experimental Setup}

\subsection{Models and Splits}

The held-out evaluation uses a stratified test split with 20 problems from each of the five difficulty levels for a total of 100 test items.
We evaluate paired \zs{}/\CoT{} runs for six model families:
GPT-5.2, gpt-oss-120b, Qwen~3.5 Flash, Gemini~3 Flash, DeepSeek-v3.2, and Llama~3.3~70B.
Additional exploratory runs exist for other models, but they are not part of the main comparison reported here.

\subsection{Evaluation Harness}

The evaluation harness executes a prompt, parses the returned JSON, computes metrics, and attaches an error report.
For non-\zs{} conditions, the harness retries once with a \zs{} prompt if the initial parse fails.
This matters empirically: some of the apparent \CoT{} gains, particularly for gpt-oss-120b and Llama~3.3~70B, are partly format-recovery effects rather than pure reasoning gains.

All runs use deterministic decoding settings with temperature fixed at \(0\).
A rule-based baseline provides an oracle-style upper bound and reaches \fbs{} $=99.97$ on the held-out split.

\subsection{Secondary Analysis Artifacts}

Three retained result summaries are especially important for the paper: primary leaderboard metrics, a deep analysis bundle covering complexity attribution and difficulty curves, and a significance bundle with bootstrap confidence intervals and paired permutation tests.
The significance bundle does not cover every model pair, so the paper treats some comparisons as descriptive rather than inferential.
Where exact \(p\)-values are cited below, they come directly from the retained permutation-test output.
